% GENERATED FILE. Edit canonical lessons, then run npm run generate:books.
% canonical-source-hash: fnv1a64:5ea79d13a231935b
% canonical-lessons: RU-C306-pomolvka, RU-C306-vlyublyatsya, RU-C306-ukhazhivat, RU-C306-venchaniye, RU-C306-rasstavatsya

\chapter{Love and engagement}
\label{ch:ru-love-and-engagement}
\hlchaptermodality{306}

\begin{chapteropening}
\textbf{By the end of this chapter:} \emph{I can talk about an engagement, falling in love, courting, a church wedding and breaking up.}

Complete the last lesson of chapter 306: I can talk about an engagement, falling in love, courting, a church wedding and breaking up.
\end{chapteropening}

\section[\texorpdfstring{\ru{помолвка} (pomólvka)}{помолвка (pomólvka)}]{\ru{помолвка} (pomólvka) — an engagement}

\label{lesson:RU-C306-pomolvka}



\begin{quote}
Before the new one: say the Russian for a graduation party, then the Russian for a generation.
\end{quote}

\subsection*{You'll want to know: \ru{помолвка}}
\textbf{\ru{помолвка}} (\emph{pomólvka}) — \textquotedblleft{}an engagement\textquotedblright{}.

\begin{quote}
\ru{После} \ru{помолвки} \ru{они} \ru{поехали} \ru{к} \ru{родителям}.
\end{quote}

\begin{quote}
\emph{Pósle pomólvki oní poyékhali k rodítelyam.}
\end{quote}

\begin{quote}
After the engagement they went to visit the parents.
\end{quote}

Feminine.

\begin{cousinweb}
\emph{\ru{помолвка}} is built on \emph{\ru{помолвить}} \textquotedblleft{}to promise in marriage\textquotedblright{}, from \emph{\ru{молвить}} \textquotedblleft{}to say\textquotedblright{}, inherited from Proto-Slavic \emph{mŭlviti} (EtymDB): a spoken promise.
\end{cousinweb}

\begin{grammarlens}[title={What you've built}]
One more word for telling, explaining and describing at B1.
\end{grammarlens}

\subsection*{Your turn}
\begin{itemize}
\raggedright
  \item \emph{Say it:} \emph{pomólvka}
  \item \emph{Say it:} \emph{pomólvka}, once more
  \item \emph{Say it:} a sentence of your own with \emph{pomólvka}
  \item \emph{Recall:} say the Russian for a graduation party, then the Russian for a generation, then say \emph{pomólvka} again
\end{itemize}

\begin{tcolorbox}[breakable,colback=teal!4,colframe=teal!35!black,title={Before you move on}]
What does \ru{помолвка} mean? (\textquotedblleft{}An engagement\textquotedblright{}.) Say it once more.
\end{tcolorbox}
\section[\texorpdfstring{\ru{влюбляться} (vlyublyátsya)}{влюбляться (vlyublyátsya)}]{\ru{влюбляться} (vlyublyátsya) — to fall in love}

\label{lesson:RU-C306-vlyublyatsya}



\begin{quote}
Before the new one: say the Russian for a generation, then the Russian for an engagement.
\end{quote}

\subsection*{You'll want to know: \ru{влюбляться}}
\textbf{\ru{влюбляться}} (\emph{vlyublyátsya}) — \textquotedblleft{}to fall in love\textquotedblright{}.

\begin{quote}
\ru{Он} \ru{влюбился} \ru{в} \ru{неё} \ru{с} \ru{первого} \ru{взгляда}.
\end{quote}

\begin{quote}
\emph{On vlyubílsya v neyó s pérvovo vzglyáda.}
\end{quote}

\begin{quote}
He fell in love with her at first sight.
\end{quote}

Imperfective; perfective \emph{\ru{влюбиться}} (\emph{vlyubítsya}), \emph{\ru{влюблюсь}}. It takes \emph{\ru{в}} + accusative.

\begin{cousinweb}
\emph{\ru{влюбляться}} is built on \emph{\ru{в}-} \textquotedblleft{}into\textquotedblright{} and \emph{\ru{любить}} \textquotedblleft{}to love\textquotedblright{}, inherited from Proto-Slavic \emph{ľubiti}, from PIE \emph{lewb\textsuperscript{h}-} (EtymDB); English \emph{love} is a cousin.
\end{cousinweb}

\begin{grammarlens}[title={What you've built}]
One more word for telling, explaining and describing at B1.
\end{grammarlens}

\subsection*{Your turn}
\begin{itemize}
\raggedright
  \item \emph{Say it:} \emph{vlyublyátsya}
  \item \emph{Say it:} \emph{vlyublyátsya}, once more
  \item \emph{Say it:} a sentence of your own with \emph{vlyublyátsya}
  \item \emph{Recall:} say the Russian for a generation, then the Russian for an engagement, then say \emph{vlyublyátsya} again
\end{itemize}

\begin{tcolorbox}[breakable,colback=teal!4,colframe=teal!35!black,title={Before you move on}]
What does \ru{влюбляться} mean? (\textquotedblleft{}To fall in love\textquotedblright{}.) Say it once more.
\end{tcolorbox}
\section[\texorpdfstring{\ru{ухаживать} (ukházhivat)}{ухаживать (ukházhivat)}]{\ru{ухаживать} (ukházhivat) — to court; to look after}

\label{lesson:RU-C306-ukhazhivat}



\begin{quote}
Before the new one: say the Russian for an engagement, then the Russian for to fall in love.
\end{quote}

\subsection*{You'll want to know: \ru{ухаживать}}
\textbf{\ru{ухаживать}} (\emph{ukházhivat}) — \textquotedblleft{}to court; to look after\textquotedblright{}.

\begin{quote}
\ru{Он} \ru{красиво} \ru{ухаживал} \ru{за} \ru{ней}.
\end{quote}

\begin{quote}
\emph{On krasívo ukházhival za ney.}
\end{quote}

\begin{quote}
He courted her beautifully.
\end{quote}

Imperfective; it takes \emph{\ru{за}} + instrumental. Also to look after someone: \emph{\ru{ухаживать} \ru{за} \ru{бабушкой}}.

\begin{cousinweb}
\emph{\ru{ухаживать}} is built on \emph{\ru{у}-} and \emph{\ru{хаживать}}, an old repeated form of \emph{\ru{ходить}} \textquotedblleft{}to go\textquotedblright{}, inherited from Proto-Slavic \emph{xoditi} (EtymDB): to keep going round someone.
\end{cousinweb}

\begin{grammarlens}[title={What you've built}]
One more word for telling, explaining and describing at B1.
\end{grammarlens}

\subsection*{Your turn}
\begin{itemize}
\raggedright
  \item \emph{Say it:} \emph{ukházhivat}
  \item \emph{Say it:} \emph{ukházhivat}, once more
  \item \emph{Say it:} a sentence of your own with \emph{ukházhivat}
  \item \emph{Recall:} say the Russian for an engagement, then the Russian for to fall in love, then say \emph{ukházhivat} again
\end{itemize}

\begin{tcolorbox}[breakable,colback=teal!4,colframe=teal!35!black,title={Before you move on}]
What does \ru{ухаживать} mean? (\textquotedblleft{}To court; to look after\textquotedblright{}.) Say it once more.
\end{tcolorbox}
\section[\texorpdfstring{\ru{венчание} (venchániye)}{венчание (venchániye)}]{\ru{венчание} (venchániye) — a church wedding}

\label{lesson:RU-C306-venchaniye}



\begin{quote}
Before the new one: say the Russian for to fall in love, then the Russian for to court.
\end{quote}

\subsection*{You'll want to know: \ru{венчание}}
\textbf{\ru{венчание}} (\emph{venchániye}) — \textquotedblleft{}a church wedding\textquotedblright{}.

\begin{quote}
\ru{Венчание} \ru{было} \ru{в} \ru{старой} \ru{церкви}.
\end{quote}

\begin{quote}
\emph{Venchániye býlo v stároy tsérkvi.}
\end{quote}

\begin{quote}
The church wedding was in an old church.
\end{quote}

Neuter. In the Orthodox rite the couple are crowned.

\begin{cousinweb}
\emph{\ru{венчание}} is built on \emph{\ru{венчать}} \textquotedblleft{}to crown, to wed\textquotedblright{}, from \emph{\ru{венец}} \textquotedblleft{}a crown, a wreath\textquotedblright{} (EtymDB lists \emph{\ru{венец}} with no parent): crowns are held over the couple.
\end{cousinweb}

\begin{grammarlens}[title={What you've built}]
One more word for telling, explaining and describing at B1.
\end{grammarlens}

\subsection*{Your turn}
\begin{itemize}
\raggedright
  \item \emph{Say it:} \emph{venchániye}
  \item \emph{Say it:} \emph{venchániye}, once more
  \item \emph{Say it:} a sentence of your own with \emph{venchániye}
  \item \emph{Recall:} say the Russian for to fall in love, then the Russian for to court, then say \emph{venchániye} again
\end{itemize}

\begin{tcolorbox}[breakable,colback=teal!4,colframe=teal!35!black,title={Before you move on}]
What does \ru{венчание} mean? (\textquotedblleft{}A church wedding\textquotedblright{}.) Say it once more.
\end{tcolorbox}
\section[\texorpdfstring{\ru{расставаться} (rasstavátsya)}{расставаться (rasstavátsya)}]{\ru{расставаться} (rasstavátsya) — to part, to break up}

\label{lesson:RU-C306-rasstavatsya}



\begin{quote}
Before the new one: say the Russian for to court, then the Russian for a church wedding.
\end{quote}

\subsection*{You'll want to know: \ru{расставаться}}
\textbf{\ru{расставаться}} (\emph{rasstavátsya}) — \textquotedblleft{}to part, to break up\textquotedblright{}.

\begin{quote}
\ru{Они} \ru{расстались} \ru{друзьями}.
\end{quote}

\begin{quote}
\emph{Oní rasstális druzyámi.}
\end{quote}

\begin{quote}
They parted as friends.
\end{quote}

Imperfective; perfective \emph{\ru{расстаться}} (\emph{rasstátsya}), \emph{\ru{расстанусь}}. Present \emph{\ru{расстаюсь}}, \emph{\ru{расстаёшься}} (\emph{rasstayóshsya}). It takes \emph{\ru{с}} + instrumental.

\begin{cousinweb}
\emph{\ru{расставаться}} is built on \emph{\ru{рас}-} \textquotedblleft{}apart\textquotedblright{} and \emph{\ru{стать}} \textquotedblleft{}to stand\textquotedblright{}, inherited from Proto-Slavic \emph{stati}, from PIE \emph{steh\textsubscript{2}-} \textquotedblleft{}to stand\textquotedblright{} (EtymDB): to stand apart.
\end{cousinweb}

\begin{grammarlens}[title={What you've built}]
One more word for telling, explaining and describing at B1.
\end{grammarlens}

\subsection*{Your turn}
\begin{itemize}
\raggedright
  \item \emph{Say it:} \emph{rasstavátsya}
  \item \emph{Say it:} \emph{rasstavátsya}, once more
  \item \emph{Say it:} a sentence of your own with \emph{rasstavátsya}
  \item \emph{Recall:} say the Russian for to court, then the Russian for a church wedding, then say \emph{rasstavátsya} again
\end{itemize}

\begin{tcolorbox}[breakable,colback=teal!4,colframe=teal!35!black,title={Before you move on}]
What does \ru{расставаться} mean? (\textquotedblleft{}To part, to break up\textquotedblright{}.) Say it once more.
\end{tcolorbox}
